\section{The joint electron--lattice response}\label{app:D.33}

We derive the properties of the joint response \cref{eq:12.1-joint} and the Migdal entry \cref{eq:12.1-Skk}. The notation is that of \cref{sec:8.6,app:D.28}: adiabatic electronic states $\Psi_j(u)$ with energies $E_j(u)$, $\Delta E_j = E_j - E_0$, the derivative coupling $\vect{d}^I_{j0} = \langle\Psi_j|\nabla_{\vect{X}_I}\Psi_0\rangle$ and the force matrix element $\vect{F}^I_{j0} = \langle\Psi_j|\nabla_{\vect{X}_I}V_{eN}|\Psi_0\rangle$.

\subsection{Labels, marginal and positivity}\label{app:D.33-labels}

Let $H_{\mathrm{el}} = \sum_\Lambda E^{\mathrm{el}}_\Lambda P_\Lambda$ be the operator that multiplies each final-state class $\Lambda$ of \cref{sec:8.7-classes} by its electronic energy, and $H_{\mathrm{lat}} = H - H_{\mathrm{el}}$. Where the classes are good labels, $[P_\Lambda, H] = 0$, and then $H_{\mathrm{el}}$ and $H_{\mathrm{lat}}$ commute with $H$ and with each other. Their joint eigenstates are the eigenstates $|f\rangle$ of $H$, with eigenvalues $E^{\mathrm{el}}_f$ and $E_f - E^{\mathrm{el}}_f$, and \cref{eq:12.1-joint} can be written without reference to a basis,
\begin{equation*}
  S_{ab}(\vect{q};\omega_l,\omega_e) = \frac{2\pi}{\Vol}\sum_ip_i\,\bigl\langle i\big|O_a^\dagger\,\delta\bigl(\omega_e + E^{\mathrm{el}}_i - H_{\mathrm{el}}\bigr)\,\delta\bigl(\omega_l + E_i - E^{\mathrm{el}}_i - H_{\mathrm{lat}}\bigr)\,O_b\big|i\bigr\rangle .
\end{equation*}
Integrating over $\omega_e$ at fixed $\omega = \omega_l + \omega_e$ turns the product of the two delta functions into $\delta(\omega + E_i - H_{\mathrm{el}} - H_{\mathrm{lat}}) = \delta(\omega + E_i - H)$, since $H_{\mathrm{el}}$ and $H_{\mathrm{lat}}$ commute, and leaves \cref{eq:4.2-Sab}, which is the marginal \cref{eq:12.1-marginal}. For any complex vector $c_a$, the contraction $\sum_{ab}c_a^*S_{ab}c_b$ at a point of the plane is $(2\pi/\Vol)\sum_{i,f}p_i|\langle f|\sum_bc_bO_b|i\rangle|^2$ times the two delta functions, a sum of nonnegative terms. The matrix $S_{ab}$ is therefore positive semidefinite at every $(\omega_l, \omega_e)$, and the bound $|S_{e\kappa}|^2 \le S_{ee}S_{\kappa\kappa}$ of \cref{sec:12.1-def} is the positivity of its two-by-two minor. Where the labels commute with $H$ only to the accuracy of \cref{eq:8.6-leak} or of the off-shell window, the marginal and the positivity hold on bins of the plane, with the remainders of \cref{eq:8.6-decomp,eq:8.6-bin}.

At $T = 0$ the initial state is the vibrational ground state of the ground surface, with $E^{\mathrm{el}}_i = 0$. A final state in $\mathcal{P}_{\mathrm{ad}}$ has $E^{\mathrm{el}}_f = 0$, which places the electron-following sector on the line $\omega_e = 0$. A final state in $\mathcal{P}_0$ has the zero-phonon energy of its excitation as its total excitation energy, $E_f - E_i = E^{\mathrm{el}}_f$, which places the final states of $\mathcal{P}_0$ on the line $\omega_l = 0$. Every other final state has $\omega_e > 0$ and $\omega_l > 0$. The placement of the three kinds of final state in the plane is \cref{eq:12.1-sectors}.

\subsection{The force operator and the amplitude}\label{app:D.33-amp}

Let the electrons feel nucleus $I$, of species $\kappa$, through a potential that moves rigidly with it, $V_{eN} \ni \sum_ev_\kappa(\vect{r}_e - \vect{X}_I)$. With the transform \cref{eq:A.3-space}, $v_\kappa(\vect{r}) = \int\dd^3p\,(2\pi)^{-3}\ee^{-\ii\vect{p}\cdot\vect{r}}v_\kappa(\vect{p})$, and with $v_\kappa(-\vect{p}) = v_\kappa(\vect{p})$ for a spherical potential,
\begin{equation*}
  \sum_ev_\kappa(\vect{r}_e - \vect{X}_I) = \int\frac{\dd^3p}{(2\pi)^3}\;v_\kappa(\vect{p})\,\ee^{-\ii\vect{p}\cdot\vect{X}_I}\sum_e\ee^{\ii\vect{p}\cdot\vect{r}_e}
  = \int\frac{\dd^3p}{(2\pi)^3}\;v_\kappa(\vect{p})\,\ee^{-\ii\vect{p}\cdot\vect{X}_I}\,n_e(\vect{p}) ,
\end{equation*}
after the substitution $\vect{p} \to -\vect{p}$, with $n_e(\vect{p})$ of \cref{eq:A.3-ne}. Only the term $\sum_ev_\kappa(\vect{r}_e - \vect{X}_I)$ of $V_{eN}$ depends on $\vect{X}_I$. Its gradient gives \cref{eq:12.1-force}, and its matrix element between adiabatic states gives
\begin{equation}
  \vect{q}\cdot\vect{F}^I_{j0} = -\ii\int\frac{\dd^3p}{(2\pi)^3}\;(\vect{q}\cdot\vect{p})\,v_\kappa(\vect{p})\,\bigl\langle\Psi_j(u)\big|\,n_e(\vect{p})\,\big|\Psi_0(u)\bigr\rangle\,\ee^{-\ii\vect{p}\cdot\vect{X}_I} .
  \label{eq:D.33-F}
\end{equation}
Both $\ee^{-\ii\vect{p}\cdot\vect{X}_I}$ and the transition density depend on the nuclear positions, and the product is a function of $u$ that multiplies the nuclear wavefunction.

The boost amplitude \cref{eq:8.6-boost} is $-\ii\,\vect{q}\cdot\vect{d}^I_{j0}/(M_\kappa\Delta E_j)$ times $\ee^{\ii\vect{q}\cdot\vect{X}_I}$ between the vibrational states. By \cref{eq:5.1-dcoup}, $\vect{d}^I_{j0} = \vect{F}^I_{j0}/(E_0 - E_j) = -\vect{F}^I_{j0}/\Delta E_j$, so the electronic factor is $\ii\,\vect{q}\cdot\vect{F}^I_{j0}/(M_\kappa\Delta E_j^2)$. The derivation of \cref{eq:D.28-boost} treats the derivative coupling as a constant over the vibrational wavefunctions. Keeping the derivative coupling inside the nuclear matrix element, as a function of $u$, gives the same leading term together with the Herzberg--Teller corrections. With \cref{eq:D.33-F} and $\ii\cdot(-\ii) = 1$,
\begin{equation*}
  A_{fj} = \sum_{I\in\kappa}\int\frac{\dd^3p}{(2\pi)^3}\;\frac{\vect{q}\cdot\vect{p}}{M_\kappa\Delta E_j^2}\,v_\kappa(\vect{p})\,
  \bigl\langle\Psi^{\mathrm{nuc}\prime}\big|\,\langle\Psi_j(u)|n_e(\vect{p})|\Psi_0(u)\rangle\,\ee^{\ii(\vect{q} - \vect{p})\cdot\vect{X}_I}\big|\Psi^{\mathrm{nuc}}\bigr\rangle .
\end{equation*}
The Condon approximation in the frame of the crystal replaces the transition density by its value at $u = 0$, which leaves the operator $\ee^{\ii(\vect{q}-\vect{p})\cdot\vect{X}_I}$ between the vibrational states and gives \cref{eq:12.1-amp}. The first correction to the Condon approximation is the gradient of the transition density with respect to the displacements, the Herzberg--Teller term of \cref{eq:8.6-electronic}.

If instead the electrons move rigidly with nucleus $I$, as in an isolated atom, the transition density at nuclear position $\vect{X}_I$ is the transition density at the origin times $\ee^{\ii\vect{p}\cdot\vect{X}_I}$, since translating every electron by $\vect{X}_I$ multiplies $n_e(\vect{p})$ by this phase. The phase cancels the factor $\ee^{-\ii\vect{p}\cdot\vect{X}_I}$ of \cref{eq:D.33-F}, and the operator between the vibrational states becomes $\ee^{\ii\vect{q}\cdot\vect{X}_I}$, a recoil with the full momentum $\vect{q}$. The frame of the crystal and the frame of the nucleus differ by the dependence of the transition density on $u$, which the frame of the crystal neglects and the frame of the nucleus overstates. Expanded in $\vect{p}\cdot\vect{u}_I$, the lattice factors of the two frames agree at zeroth order and differ at first order: the lattice keeps $\vect{q} - \vect{p}$ in the frame of the crystal and $\vect{q}$ in the frame of the nucleus, which matters for the lattice spectrum when $p$ is not small against $q$ (\cref{sec:12.4}).

\subsection{Squared amplitude and crystal momentum}\label{app:D.33-square}

In a box of volume $\Vol$, $\int\dd^3p\,(2\pi)^{-3} \to \Vol^{-1}\sum_{\vect{p}}$. The squared amplitude summed over the nuclei of species $\kappa$ is
\begin{equation*}
  \begin{split}
  |A_{fj}|^2 = \frac{1}{M_\kappa^2\Delta E_j^4\Vol^2}\sum_{\vect{p}\vect{p}'}\;&(\vect{q}\cdot\vect{p})(\vect{q}\cdot\vect{p}')\,v_\kappa(\vect{p})^*v_\kappa(\vect{p}')\,
  \langle\Psi_0|n_e(\vect{p})^\dagger|\Psi_j\rangle\langle\Psi_j|n_e(\vect{p}')|\Psi_0\rangle \\
  &\times\langle\Psi^{\mathrm{nuc}}|n_\kappa(\vect{q}-\vect{p})^\dagger|\Psi^{\mathrm{nuc}\prime}\rangle\langle\Psi^{\mathrm{nuc}\prime}|n_\kappa(\vect{q}-\vect{p}')|\Psi^{\mathrm{nuc}}\rangle ,
  \end{split}
\end{equation*}
with $n_\kappa(\vect{Q}) = \sum_{I\in\kappa}\ee^{\ii\vect{Q}\cdot\vect{X}_I}$, the transition densities at $u = 0$, and the real factor $\vect{q}\cdot\vect{p}$. The final state has $\omega_e = \Delta E_j$ and $\omega_l$ equal to the energy of $\Psi^{\mathrm{nuc}\prime}$ above $\Psi^{\mathrm{nuc}}$. After multiplication by $(2\pi/\Vol)p_i\,\delta(\omega_e - \Delta E_j)\,\delta(\omega_l - E^{\mathrm{nuc}\prime} + E^{\mathrm{nuc}})$ and summation over $j \neq 0$, over $\Psi^{\mathrm{nuc}\prime}$ and over the thermal initial states, the electronic sum is $(\Vol/2\pi)$ times the density response between $n_e(\vect{p})$ and $n_e(\vect{p}')$ at clamped nuclei, and the nuclear sum is $(\Vol/2\pi)$ times the electron-following response between $n_\kappa(\vect{q}-\vect{p})$ and $n_\kappa(\vect{q}-\vect{p}')$. In a crystal the clamped electrons conserve crystal momentum, and their response vanishes unless $\vect{p}$ and $\vect{p}'$ differ by a reciprocal-lattice vector. With $\vect{p} = \vect{k} + \vect{G}$ and $\vect{p}' = \vect{k} + \vect{G}'$ and $\vect{k}$ in the zone, the double sum becomes $\sum_{\vect{k}}\sum_{\vect{G}\vect{G}'}$, and $\Vol^{-1}\sum_{\vect{k}} \to \int_{\mathrm{BZ}}\dd^3k\,(2\pi)^{-3}$. Collecting the factors gives $(2\pi/\Vol)(\Vol/2\pi)^2\Vol^{-2}\cdot \Vol = (2\pi)^{-1}$, and with $\Delta E_j = \omega_e$ under the delta function we obtain \cref{eq:12.1-Skk}.
